
% !TEX root = main.tex

\subparagraph*{Shifting}

Given a grade $\sgr$ and a pair of grades $\ple{\Rgr,\rgr}$, where the first one is regarded as an assumption grade, we define the \emph{shif} of $\ple{\Rgr,\rgr}$ by $\sgr$ as the operation 
\begin{quote}
$\sgr\rsh\ple{\Rgr,\rgr} = \ple{\Rgr\rsum\sgr\rmul\rgr,\rgr}$. 
\end{quote}
This operation extends to types and contexts in the obvious way. 

\begin{proposition}\label{prop:rshift}
If $\IsWFField\Sigma\mem\val\T$, then, for every grade $\sgr$ we have 
$\IsWFField{\sgr\rsh\Sigma}{\mem}{\val}{\sgr\rsh\T}$. 
\end{proposition}
\begin{proof}
The proof is by coinduction. 
The only non-trivial case is when $\val = \loc$. 
Consider a judgement 
$\IsWFField{\sgr\rsh\Sigma}{\mem}{\loc}{\Graded\C{\Rgr\rsum\sgr\rmul\rgr}{\rgr}}$ such that $\IsWFField\Sigma\mem\loc{\Graded\C\Rgr\rgr}$ is derivable. 
We have to show that this is the consequence of an instance of the rule \refToRule{t-ref-in-mem} where the premises are appropriately shifted. 
Consider a judgement 
$\IsWFField{\sgr\rsh\Sigma}{\mem}{\loc}{\Graded\C{\Rgr\rsum\sgr\rmul\rgr}{\rgr}}$ such that $\IsWFField\Sigma\mem\loc{\Graded\C\Rgr\rgr}$ is derivable. 
We have to show that this is the consequence of an instance of the rule \refToRule{t-ref-in-mem} where the premises are appropriately shifted. 
By hypothesis we know that 
$\mem(\loc) = \Pair{\_}
{\Obj\C{\GradedValInd{\val}{1}{\rgr},\ldots,\GradedValInd\val{n}{\rgr}}} $, 
$\Sigma = \CtxElem{\loc}{\Rgr}{\rgr}{\C}\ctxsum \rgr\rmul\Sigma'$ and, for all $i \in 1..n$, 
$\IsWFField{\Sigma_i}{\mem}{\val_i}{\T_i}$ with 
$\T_i\subt\Graded{\tau_i}{\Rgr_i\ctxsum\Rgr\rmul\rgr_i}{\rgr_i}$, 
$\fields{\C} = \Field{\Graded{\tau_1}{\Rgr_1}{\rgr_1}}{\f_1} \ldots \Field{\Graded{\tau_n}{\Rgr_n}{\rgr_n}}{\f_n}$, 
and $\Sigma_i \ctxord \Sigma'$. 
We consider as premises the judgements 
$\IsWFField{(\sgr\rmul\rgr)\rsh\Sigma_i}{\mem}{\val_i}{(\sgr\rmul\rgr)\rsh\T_i}$  and we notice that 
$(\sgr\rmul\rgr)\rsh\Sigma_i \ctxord (\sgr\rmul\rgr)\rsh \Sigma$ and 
$(\sgr\rmul\rgr)\rsh\T_i 
  \subt (\sgr\rmul\rgr) \rsh \Graded{\tau_i}{\Rgr_i\rsum\Rgr\rmul\rgr_i}{\rgr_i} 
  = \Graded{\tau_i}{\Rgr_i \rsum \Rgr\rmul\rgr_i \rsum \sgr\rmul\rgr\rmul\rgr_i}{\rgr_i} 
  = \Graded{\tau_i}{\Rgr_i\rsum (\Rgr\rsum\sgr\rmul\rgr)\rmul\rgr_i}{\rgr_i}$, 
by monotonicity of the shifting operation. 
Notice also that 
$\sgr \rsh (\rgr\rmul\Sigma) = \rgr\rmul ((\sgr\rmul\rgr)\rsh\Sigma)$ and 
$\CtxElem{\loc}{\Rgr\rsum\sgr\rmul\rgr}{\rgr}{\C} = \sgr \rsh (\CtxElem\loc\Rgr\rgr\C)$. 
Then we get the thesis$\sgr\rsh -$ laxly distributes over addition of contexts. 
\FDComm{servirebbe ancora un po' di subsumption sui contesti per far quest'ultimo passaggio} 
\end{proof}

\PGComm{Prima: If\ $\IsWFField\Sigma\mem\val{\Graded\tau\Rgr{\rgr_1\rsum\rgr_2}}$}
\begin{proposition}\label{prop:split-ty} 
If $\IsWFField\Sigma\mem\val{\Graded\tau\Rgr{\bpa\rgr\epa}}$  \bpa and  $\rgr_1\rsum\rgr_2\rord\rgr$ \epa then 
there are $\Sigma_1$ and $\Sigma_2$ such that 
$\Sigma_1\ctxsum\Sigma_2\ctxord\Sigma$, 
$\IsWFField{\Sigma_1}{\mem}{\val}{\Graded\tau{\Rgr\rsum\rgr_2}{\rgr_1}}$ and 
$\IsWFField{\Sigma_2}{\mem}{\val}{\Graded\tau{\Rgr\rsum\rgr_1}{\rgr_2}}$. 
\end{proposition} 
\begin{proof}
The only non-trivial case is when $\val = \loc$. 
By \refToRule{t-ref-in-mem} we know that 
$\tau = \C$, 
$\mem(\loc) = \Pair{\_}
{\Obj\C{\GradedValInd{\val}{1}{\sgr},\ldots,\GradedValInd\val{n}{\sgr}}} $, 
$\Sigma = \CtxElem{\loc}{\Rgr}{\rgr}{\C}\ctxsum \rgr\rmul\Sigma'$ and, for all $i \in 1..n$, 
$\IsWFField{\Sigma'_i}{\mem}{\val_i}{\T_i}$ with 
$\T_i\subt\Graded{\tau_i}{\Sgr_i\rsum \Rgr\rmul\sgr_i}{\sgr_i}$, 
$\fields{\C} = \Field{\Graded{\tau_1}{\Sgr_1}{\sgr_1}}{\f_1} \ldots \Field{\Graded{\tau_n}{\Sgr_n}{\sgr_n}}{\f_n}$, 
and $\Sigma'_i \ctxord \Sigma'$. 
Now, by \cref{prop:rshift}, we obtain 
$\IsWFField{\rgr_2\rsh\Sigma'_i}{\mem}{\val_i}{\rgr_2\rsh\T_i}$ with 
$\rgr_2\rsh\Sigma'_i \ctxord \rgr_2\rsh\Sigma'$ and 
$\rgr_2\rsh\T_i \subt \Graded\C{\Sgr_i\rsum(\Rgr\rsum\rgr_2)\rmul\sgr_i}{\sgr_i}$, 
for all $i\in 1..n$. 
Let us set 
$\Sigma_1 = \CtxElem\loc{\Rgr\rsum\rgr_2}{\rgr_1}\C \ctxsum \rgr_1\rmul (\rgr_2\rsh\Sigma')$, then, by rule \refToRule{t-ref-in-mem}, we obtain 
$\IsWFField{\Sigma_1}{\mem}{\loc}{\Graded\C{\Rgr\rsum\rgr_2}{\rgr_1}}$. 
In a similar way we derive $\IsWFField{\Sigma_2}{\mem}{\loc}{\Graded\C{\Rgr\rsum\rgr_1}{\rgr_2}}$ where $\Sigma_2  = \CtxElem\loc{\Rgr\rsum\rgr_1}{\rgr_2}\C \ctxsum \rgr_2\rmul(\rgr_1\rsh\Sigma')$. 
Therefore, to conclude, we just have to prove that $\Sigma_1\ctxsum\Sigma_2 \ctxord \Sigma$, but this is immediate by construction, as in each side we have increased the assumption grade of each reference exactly by its usage in the other side \bpa and  $\rgr_1\rsum\rgr_2\rord\rgr$. \epa

\end{proof}


